\section{Threats to Validity}
\label{sec:threats}

Several factors condition the interpretation of our results. First, to establish an un-tuned algorithmic baseline, all Boolean SAT encodings were executed under the Glucose3 default phase policy. Modifying branching heuristics could radically alter search trajectories; therefore, the timeouts and non-monotonic behavior observed in the SAT-Product encoding may be artifacts of this specific heuristic interaction rather than inherent failures of the encoding space. Second, the primary metric (pipeline total time) aggregates CNF construction and search. It is critical to distinguish between per-group complexity and total formula complexity; while some encodings have $O(N^2)$ clauses per AMO group yielding a total of $\Theta(N^3)$ clauses (e.g., Pairwise), others require $O(N)$ clauses per group yielding $O(N^2)$ total clauses. Consequently, all SAT algorithms are burdened by structural initialization penalties not experienced by native $O(N)$ CP architectures.

The supplementary large-scale experiment at $N=200$ introduces additional validity considerations. Execution was limited to only five repetitions per method, which may be insufficient to capture the full variance of highly stochastic CDCL search paths. Furthermore, heavy swap usage and strict preflight memory exclusions (\texttt{NOT\_RUN\_RESOURCE\_POLICY}) were necessitated by the hardware constraints of the single macOS ARM64 environment. The IBM CP Optimizer's \texttt{LICENSE\_ERROR} was explicitly triggered by a commercial $2^{1000}$ search-space bound rather than algorithmic inability. Finally, adding a single large $N=200$ point provides valuable extreme-scale insights but is not sufficient to definitively project continuous asymptotic complexity curves across all intermediate problem sizes.
